\documentclass[10pt,a4paper]{amsart}
\usepackage[margin=19mm]{geometry}
\usepackage{amsmath,amssymb,mathtools}
\usepackage[scr=rsfs,cal=euler]{mathalfa}
\usepackage{xfrac}
\usepackage[unicode,hidelinks,bookmarks=false]{hyperref}
\newcommand{\ie}{i.e.\ }
\newcommand{\cf}{cf.\ }
\newcommand{\resp}{respectively\ }
\newcommand{\Subsection}{\subsection}
\providecommand{\nobreakdash}{\nobreak\hbox{-}}
\setlength{\emergencystretch}{2em}
\multlinegap0pt
\usepackage{amssymb}
\usepackage{mathtools}
\usepackage{float}
\usepackage[shortlabels]{enumitem}
\newtheorem{lemma}{Lemma}
\newtheorem{theorem}{Theorem}
\newtheorem{remark}{Remark}
\DeclareMathOperator{\Li}{Li}
\title{${}_5F_4$ evaluations and a family of $\pi^2+\log^2$ identities}
\author{Cetin Hakimoglu-Brown}
\date{Source: arXiv:2608.19276v1 (CC BY 4.0)}
\begin{document}
\maketitle

\noindent\emph{Source and licence.} ${}_5F_4$ evaluations and a family of $\pi^2+\log^2$ identities. Version arXiv:2608.19276v1 (CC BY 4.0), distributed by arXiv under the Creative Commons Attribution 4.0 International licence, http://creativecommons.org/licenses/by/4.0/, as recorded in the arXiv licence metadata for this exact version and shown on the abstract page https://arxiv.org/abs/2608.19276v1. Full licence notice: \texttt{licenses/arxiv-2608.19276-CC-BY-4.0.txt}.\par\smallskip

\noindent\textit{Editorial scope.} Counted unit: the article's lemma lem:landen (source0.tex lines 202--209). The article's complete original proof of it is delivered with it (source0.tex lines 211--247, one \texttt{proof} environment of 37 lines). No mathematics is written, repaired or re-derived: the statement and the proof are the article's own lines, in the article's own notation, and no retained line is substituted or re-flowed.  Retained uncounted as the source-local context the proof needs: lines 122--124; lines 315--322; lines 369--374; lines 401--417; lines 419--441.  Nothing was omitted from any retained span, so the package carries no omission record.  No retained line contains a citation, so the wrapper carries no bibliography and no bibliography entry.  No retained line contains a figure environment, an image-inclusion command or drawing code, so no image is delivered and the package declares no asset dependency.  Editorial formatting only, in addition: an AMS \texttt{amsart} preamble that restates the article's own declarations and macros used by the retained lines, namely the environments \texttt{lemma}, \texttt{theorem}, \texttt{remark} on separate counters, together with the standard packages those lines need.  The retained environments print in this document as Lemma 1, Theorem 1, Remark 1 in order of appearance, whereas the article prints the same items with the numbering of its own full document.  The complete native source of the article as distributed in the arXiv e-print for this exact version is bundled unchanged as \texttt{sources/208075/source0.tex}.
\begin{equation}\label{eq:fiber}
  x^4-3x^3+3x^2-x+\tfrac1z=0,\qquad\text{equivalently}\qquad x(1-x)^3=\tfrac1z .
\end{equation}

\begin{lemma}\label{lem:angle}
For an admissible member with forced angle $\theta=j\pi/N$ and $z<0$,
\[
  \arg\!\big(1-\tfrac1\rho\big)=\tfrac23(\pi-2\theta),
  \qquad\text{hence}\qquad
  -\tfrac38\cdot 2\arg^2\!\big(1-\tfrac1\rho\big)=-\tfrac13\Big(1-\tfrac{2j}{N}\Big)^2\pi^2 .
\]
\end{lemma}

\begin{lemma}\label{lem:modulus}
For an admissible member with forced angle $\theta$ one has $\theta<\pi/2$ and
\[
  r\;=\;\frac{\sin\!\big(\tfrac{\pi+\theta}{3}\big)}{\sin\!\big(\tfrac{\pi+4\theta}{3}\big)} .
\]
\end{lemma}

\begin{theorem}\label{thm:main}
Let $\theta=j\pi/N$ with $\gcd(j,N)=1$, and let $r$ be a positive real root of
$(\kappa^3-2\kappa)r^3+3(1-\kappa^2)r^2+3\kappa r-1=0$ with $\kappa=2\cos\theta$; for an admissible
member this is the value $r=\sin\big(\tfrac{\pi+\theta}{3}\big)/\sin\big(\tfrac{\pi+4\theta}{3}\big)$
of Lemma~\textup{\ref{lem:modulus}}. Put
$S=3-\kappa r$, $P=3-r^2-3\kappa r+\kappa^2r^2$, $z=1/(Pr^2)$ and
$a,b=\tfrac12\big(S\pm\sqrt{S^2-4P}\big)$, with $a$ the larger root. If $S^2-4P>0$ and $|z|<256/27$,
then necessarily $z<0$ and $a>0>b$ (Remark~\textup{\ref{rem:signs}}), the remaining two roots of
\eqref{eq:fiber} are $re^{\pm i\theta}$, and
\begin{align*}
  -\frac{z}{4}\;{}_5F_4\!\left(\begin{matrix}1,1,1,\tfrac43,\tfrac53\\[2pt]
      \tfrac54,\tfrac32,\tfrac74,2\end{matrix}\;\middle|\;\frac{27z}{256}\right)
  &= -\frac13\Big(1-\frac{2j}{N}\Big)^2\pi^2
    + \tfrac12\ln^2 a+\tfrac12\ln^2|b|+\tfrac23\ln^2 r\\[2pt]
  &\qquad-\tfrac13\ln a\ln|b|-\tfrac23\ln r\,(\ln a+\ln|b|).
\end{align*}
\end{theorem}

\begin{remark}[the signs of $z$, $a$, $b$ are automatic]\label{rem:signs}
An admissible member has $z<0$. Indeed $x(1-x)^3\le 27/256$ on $[0,1]$ and is negative outside it, so
a real root of $x(1-x)^3=1/z$ with $z>0$ requires $1/z\le27/256$, i.e.\ $z\ge z_c=256/27$. Hence any
member with real $a,b$ and $|z|<z_c$ has $z<0$, which is the hypothesis used in Lemma~\ref{lem:angle}.
The cubic in $r$ may have more than one positive root; only a branch with $S^2-4P>0$ and $|z|<z_c$ is
admissible, and it is this branch that is meant throughout.

The signs of $a$ and $b$ are likewise automatic, so that $\ln a$ and $\ln|b|$ in
Theorem~\ref{thm:main} are the real logarithms of positive reals. Indeed the product of the four
roots equals the constant term $1/z$ of \eqref{eq:fiber}, and the conjugate pair contributes
$|\rho|^2=r^2>0$; hence $ab\,r^2=1/z<0$, so $ab<0$ and $a,b$ have opposite signs. With $a$ the larger
root this gives
\[
  a>0>b ,
\]
and all logarithms occurring in Theorem~\ref{thm:main} are taken as real logarithms of the positive
quantities $a$, $|b|$ and $r$; no complex branch of $\log$ is needed. Since $a>1$ and $b<0$ we have
in addition
\[
  1-\frac1a>0,\qquad 1-\frac1b>0,
\]
which is used in the proof of Lemma~\ref{lem:landen} to control principal arguments.
\end{remark}

\begin{lemma}\label{lem:landen}
For any $z$ in the domain of convergence,
\[
  \int_0^1\frac{\ln\!\big(1-z\,x(1-x)^3\big)}{x(1-x)}\,dx
  =\frac12\sum_{k=1}^4\ln^2\!\Big(1-\frac1{\rho_k}\Big),
\]
where $\rho_1,\dots,\rho_4$ are the roots of the fiber polynomial \eqref{eq:fiber}.
\end{lemma}

\begin{proof}
For $0<\varepsilon<1$ set
\[
  I_\varepsilon=\int_0^{1-\varepsilon}\frac{\ln\!\big(1-z\,x(1-x)^3\big)}{x(1-x)}\,dx .
\]
On $[0,1-\varepsilon]$ the integrand is bounded, so the factorisation
$\ln(1-z\,x(1-x)^3)=\sum_k\ln(1-x/\rho_k)$ and the splitting
$\tfrac1{x(1-x)}=\tfrac1x+\tfrac1{1-x}$ may be performed term by term, each resulting integral being
finite:
\[
  I_\varepsilon=\sum_{k=1}^4\left[\int_0^{1-\varepsilon}\frac{\ln(1-x/\rho_k)}{x}\,dx
  +\int_0^{1-\varepsilon}\frac{\ln(1-x/\rho_k)}{1-x}\,dx\right].
\]
As $\varepsilon\to0^+$ the first integral converges to $-\Li_2(1/\rho_k)$. The second diverges
individually; integrating by parts and expanding,
\[
  \int_0^{1-\varepsilon}\frac{\ln(1-x/\rho_k)}{1-x}\,dx
  =-\ln\varepsilon\cdot\ln\!\Big(1-\frac1{\rho_k}\Big)-\Li_2\!\Big(\frac{1}{1-\rho_k}\Big)+o(1).
\]
Summing over $k$, the divergent contributions carry the total coefficient
$\sum_{k}\ln(1-1/\rho_k)$, which vanishes. Indeed $\prod_k(1-1/\rho_k)=p(1)/p(0)=1$, $p$ being the
monic fiber polynomial \eqref{eq:fiber}; and passing from $\sum_k\ln w_k$ to $\ln\prod_k w_k$ with
principal logarithms is legitimate here because $\sum_{k=1}^4\arg(1-1/\rho_k)=0$: the nonreal roots
form a conjugate pair, so $\arg(1-1/\rho)+\arg(1-1/\bar\rho)=0$, while by Remark~\ref{rem:signs}
$1-1/a>0$ and $1-1/b>0$ each contribute argument $0$. No multiple of $2\pi i$ intervenes, so
\[
  \sum_{k=1}^4\ln\Big(1-\frac1{\rho_k}\Big)=\ln 1=0 .
\]
Therefore $I_\varepsilon$ converges as $\varepsilon\to0^+$, with
\[
  \lim_{\varepsilon\to0^+}I_\varepsilon
  =-\sum_{k=1}^4\left[\Li_2\!\Big(\frac1{\rho_k}\Big)+\Li_2\!\Big(\frac{1}{1-\rho_k}\Big)\right].
\]
Finally, putting $w=1/\rho_k$ one has $\tfrac{w}{w-1}=\tfrac{1}{1-\rho_k}$, so Landen's identity
$\Li_2(w)+\Li_2\!\big(\tfrac{w}{w-1}\big)=-\tfrac12\ln^2(1-w)$ collapses each bracket to
$-\tfrac12\ln^2(1-1/\rho_k)$, which is the stated result. \qedhere
\end{proof}

\end{document}
